\documentclass[11pt]{article}
\usepackage[margin=0.68in]{geometry}
\usepackage{lmodern}
\usepackage{microtype}
\usepackage{enumitem}
\usepackage[table]{xcolor}
\usepackage{titlesec}
\setlength{\parindent}{0pt}
\setlength{\parskip}{0.28em}
\setlist[itemize]{leftmargin=1.3em,itemsep=0.08em,topsep=0.12em}
\titleformat{\section}{\large\bfseries}{}{0pt}{}
\titlespacing*{\section}{0pt}{0.7em}{0.25em}
\pagestyle{plain}

\newenvironment{playpassage}{%
  \begin{quote}\small\setlength{\parskip}{0.08em}\setlength{\parindent}{0pt}%
}{\end{quote}}
\newcommand{\speaker}[1]{\textbf{#1:}}

\begin{document}

\begin{center}
{\Large\bfseries U1L6SP --- Text Supplement}\par
{\LARGE\bfseries Prophecy, Choice, and Consequence}\par
{\large Two scenes for close reading; the whole play remains available}
\end{center}

\fcolorbox{black}{gray!8}{%
\begin{minipage}{0.94\linewidth}\small
\textbf{How to use this supplement}

You may use \textbf{anything from the play} in your paper: a line, action, choice, image, or consequence. The passages below are not the only acceptable evidence. They are lines worth focusing on because they let you examine what the prophecy says, what Macbeth considers, and what he knows before he acts.

You do not need the internet or another source. Quote exact words when they are printed here. You may accurately paraphrase other events you remember from watching the play.
\end{minipage}}

\section*{Passage 1: Act 1, Scene 3 --- After the Witches Speak}

\textbf{Context:} The witches have greeted Macbeth as Thane of Glamis, Thane of Cawdor, and future king. Messengers then confirm that Macbeth has been named Thane of Cawdor. One prediction has come true; the prediction that he will be king remains unfulfilled.

\begin{playpassage}
\speaker{Banquo} But 'tis strange:\par
And oftentimes, to win us to our harm,\par
The instruments of darkness tell us truths,\par
Win us with honest trifles, to betray's\par
In deepest consequence.

\speaker{Macbeth} Two truths are told,\par
As happy prologues to the swelling act\par
Of the imperial theme. \ldots\par
This supernatural soliciting\par
Cannot be ill, cannot be good: if ill,\par
Why hath it given me earnest of success,\par
Commencing in a truth? I am Thane of Cawdor:\par
If good, why do I yield to that suggestion\par
Whose horrid image doth unfix my hair\par
And make my seated heart knock at my ribs,\par
Against the use of nature? Present fears\par
Are less than horrible imaginings.\par
My thought, whose murder yet is but fantastical,\par
Shakes so my single state of man\par
That function is smother'd in surmise,\par
And nothing is but what is not.

\speaker{Macbeth} If chance will have me king, why, chance may crown me,\par
Without my stir.
\end{playpassage}

\section*{Words in context}

\begin{itemize}
\item \textbf{honest trifles}: small truths;
\item \textbf{prologues}: introductions or opening parts;
\item \textbf{soliciting}: prompting or tempting;
\item \textbf{earnest}: a pledge or first sign of something promised;
\item \textbf{fantastical}: existing in imagination;
\item \textbf{stir}: action or effort.
\end{itemize}

\textbf{As you choose evidence:} Notice what the witches actually predict, what Banquo warns, what Macbeth imagines, and what Macbeth says about acting. Decide for yourself what those differences prove.

\newpage

\section*{Passage 2: Act 1, Scene 7 --- Macbeth Considers the Murder}

\textbf{Context:} Duncan is Macbeth's king, kinsman, and guest. Macbeth is considering whether to murder him. Alone, Macbeth thinks through possible consequences, his obligations to Duncan, and what is pushing him toward the deed.

\begin{playpassage}
\speaker{Macbeth} If it were done when 'tis done, then 'twere well\par
It were done quickly: if the assassination\par
Could trammel up the consequence, and catch\par
With his surcease success; that but this blow\par
Might be the be-all and the end-all here,\par
But here, upon this bank and shoal of time,\par
We'ld jump the life to come. But in these cases\par
We still have judgment here; that we but teach\par
Bloody instructions, which, being taught, return\par
To plague the inventor: this even-handed justice\par
Commends the ingredients of our poison'd chalice\par
To our own lips.

He's here in double trust;\par
First, as I am his kinsman and his subject,\par
Strong both against the deed; then, as his host,\par
Who should against his murderer shut the door,\par
Not bear the knife myself.

Besides, this Duncan\par
Hath borne his faculties so meek, hath been\par
So clear in his great office, that his virtues\par
Will plead like angels, trumpet-tongued, against\par
The deep damnation of his taking-off.

I have no spur\par
To prick the sides of my intent, but only\par
Vaulting ambition, which o'erleaps itself\par
And falls on the other.
\end{playpassage}

\section*{Words in context}

\begin{itemize}
\item \textbf{trammel up}: catch, contain, or prevent from spreading;
\item \textbf{surcease}: death;
\item \textbf{judgment here}: consequences or justice in this life;
\item \textbf{double trust}: more than one strong obligation;
\item \textbf{faculties}: powers and duties of office;
\item \textbf{taking-off}: killing;
\item \textbf{spur}: something that drives action;
\item \textbf{intent}: purpose or plan.
\end{itemize}

\section*{The rest of the play is available to you}

You may connect these lines to any later consequence or choice you remember. Possibilities include Macbeth's rule after Duncan's death, Banquo's murder, the attack on Macduff's family, the later apparitions, the branches from Birnam Wood, and Macbeth's final encounter with Macduff. These are reminders, not required evidence and not conclusions.

If you use a later event, explain how it helps test your claim. You may paraphrase the event accurately; do not invent a quotation from memory.

\end{document}
